\section{性交后反应：头痛、悲伤与不适的识别与处置}

性活动结束之后，身体并不总是立刻回到平静。有一类问题容易被忽视，因为它们往往出现在"事已毕"之后——性交后头痛、性交后情绪低落（性交后悲伤）以及性交后尿路不适。它们多数是良性的生理现象，但少数情况提示需要就医。本节按"现象—机制—处置—就医指征"逐一说明。

\subsection{性交后头痛（Post-Coital Headache）}

性交后头痛指在性兴奋或性高潮前后出现的头痛，属于"性活动相关性头痛"（sex-related headache），在国际头痛分类中被归为"原发性性活动相关性头痛"（primary headache associated with sexual activity）。

\subsubsection{两种典型表现}

- \textbf{性高潮前头痛（钝痛型）}：在性兴奋逐渐升高过程中出现，呈双侧、钝胀感，像"头颅被箍紧"，随兴奋度上升而加重，常在性高潮到达前最明显，性高潮后可逐渐缓解。多见于紧张性头痛体质者。
- \textbf{性高潮期头痛（爆发型）}：在性高潮瞬间突然出现，呈剧烈的"炸裂样"头痛，达峰时间仅数秒，持续数分钟到数小时。虽然多数为良性，但因症状与蛛网膜下腔出血、脑动脉瘤破裂极为相似，首次发生、或程度前所未有者必须首先排除器质性病变。

\subsubsection{机制}

- \textbf{血压与颅压骤升}：性高潮时心率可达 110--180 次/分，收缩压可升高 20--80 mmHg，颈部和头皮肌肉剧烈收缩，可诱发血管性与肌紧张性头痛。
- \textbf{肌肉紧张}：颈部、下颌、头皮肌肉在性活动中持续收缩，形成"发力式"头痛。
- \textbf{血管因素}：部分人存在脑血管对血流动力学变化过度敏感的倾向。

\subsubsection{处置}

- \textbf{首次发作务必就医}：尤其是爆发型、程度剧烈、伴呕吐、颈项强直、意识改变或肢体无力者，应立即急诊，行头颅 CT/MRI 及脑血管检查（CTA/MRA），排除蛛网膜下腔出血与动脉瘤。
- \textbf{良性者预防用药}：反复发作且已排除器质性病变者，可在性活动前 30--60 分钟预防性服用非甾体抗炎药（如吲哚美辛、萘普生）或 $\beta$ 受体阻滞剂（如普萘洛尔），需由神经内科医生评估后使用。
- \textbf{行为调整}：避免过度屏气用力、采取较放松的体位、充分前戏、控制兴奋上升速度，有助于减少发作。

\subsection{性交后悲伤（Post-Coital Dysphoria, PCD）}

性交后悲伤指在一段双方自愿、满意的性活动（含自慰）之后，出现的莫名悲伤、焦虑、烦躁、想哭或空虚感，可持续数分钟至数小时，甚至长达一两天。

\subsubsection{打破误解}

- \textbf{与"是否享受"无关}：PCD 可以发生在性高潮之后，也可以发生在没有高潮的性活动之后；可以发生在满意的关系中，也可以发生在自慰之后。它的出现\textbf{不代表感情出了问题}，也不代表"性冷淡"或"被强迫"。
- \textbf{并不罕见}：多项针对大学生的调查显示，约 30\%--46\% 的女性、约 20\%--41\% 的男性一生中至少经历过一次 PCD，其中约 2\%--5\% 的人反复出现。它比人们想象中常见得多。

\subsubsection{可能机制}

- \textbf{神经内分泌波动}：性高潮后催产素、催乳素、内啡肽等大量释放，随后激素水平快速回落，情绪可能随之出现"反弹性低落"；催乳素升高本身也与情绪低落、性欲抑制相关。
- \textbf{亲密后的脆弱感}：高度亲密与暴露之后，防御机制重新回升，可能触发对关系、自我价值的联想与反思。
- \textbf{心理与人际因素}：既往性创伤、对性的内疚与羞耻、依恋类型（尤其焦虑型/回避型）、关系冲突、疲劳与睡眠不足，都会增加 PCD 的发生概率。

\subsubsection{应对}

- \textbf{知情与正常化}：知道"这是一种可以被理解的现象"本身就能显著减轻痛苦——最怕的是把它误读为"我不正常"或"对方不爱我"。
- \textbf{事后安抚（aftercare）}：拥抱、依偎、轻声交谈、共同呼吸几分钟，让身体从高潮的高交感状态平稳回落；不要立即起身、看手机或背对而睡。
- \textbf{沟通}：如果反复出现，应与伴侣说明"这是我的生理与情绪反应，与你无关"，避免伴侣误解为被拒绝。
- \textbf{寻求专业帮助}：若持续存在、严重影响情绪与关系，或伴随抑郁、焦虑症状，可求助性治疗师或心理科医生。

\subsection{性交后尿路不适与尿痛（蜜月膀胱炎）}

- \textbf{表现}：性交后数小时至一两天内出现尿频、尿急、尿痛、下腹坠胀，部分人出现肉眼血尿，即所谓"蜜月膀胱炎"。
- \textbf{机制}：性交时的机械摩擦将尿道口周围的肠道细菌（以大肠埃希菌为主）推入尿道，女性尿道短（约 4 cm）且直，细菌易上行至膀胱。
- \textbf{预防}：性交前后排尿（"性交后立即排尿"是核心措施）、性交前后清洗外阴、性交后多饮水、避免使用杀精剂（可破坏阴道菌群）、润滑充分以减少摩擦。
- \textbf{处理}：出现症状应就医查尿常规与尿培养，明确后规范使用抗生素；反复发作者需排查泌尿系统结构异常，以及伴侣是否需要同步治疗。

\subsection{其他性交后不适}

- \textbf{性交后出血}：详见"女性常见性健康问题"相关章节；绝经后出血、反复出血必须就医排查宫颈病变。
- \textbf{性交后全身酸痛、疲惫}：多为体力消耗、肌肉紧张与睡眠不足所致，性活动强度接近中等强度运动，不必恐慌；持续加重的疲劳需排查贫血、甲状腺功能异常、慢性感染等。
- \textbf{性交后皮疹、瘙痒}：可能与乳胶避孕套过敏、润滑剂刺激或精液过敏有关，详见相关章节。

\begin{tcolorbox}[colback=pink!5!white,colframe=red!50!black,title=贴心小叮咛]
性交后反应多数是良性的生理现象，但它也可能是身体发出的信号。\textbf{凡符合以下任一情况，请及时就医}：首次出现的剧烈"炸裂样"头痛；伴呕吐、颈项强直、意识或肢体异常；性交后反复出血；发热伴腰痛；症状持续加重或严重影响生活。把这些反应当作"需要了解的身体信息"，而不是"羞于启齿的秘密"。
\end{tcolorbox}
